\begin{lemma}[Step 1, Subcase 2.2 at a non-PE zero divider]
Let $G,e_0,[k],T,\varphi$ be a finite loopless multigraph with exact edge
identities, distinguished uncolored edge, color set, positive-rung
certified ETT, and generating coloring.  Assume maximum property,
properness, exact one-uncoloredness, and the complete six-part lower-rung
package.  Let $H$ be a good hierarchy over a hierarchy base, with
$d<q_0$ dividers.  Assume the same-rung fewer-than-$q_0$ elementarity
hypothesis for arbitrary certified inputs and generating colorings, the
large bound $|\overline\varphi(T_n)|\geq2n+11$, the separate terminal
certificate
\[
 T\text{ is }
 \left(\bigcup_{\eta\in D^H_{n,d+1}}\Gamma^d_\eta\right)^-
 \text{-closed under }\varphi,
\]
and elementarity of $V(T)$.  Put $A=H.\mathit{starred}(d)$.

Let $\alpha\ne\beta$, with $\alpha$ missing and singleton-closed on $T$
and $\beta$ not missing on $T$.  Let $P_1,P_2,P_3$ be pairwise
support-disjoint exact $(T,\varphi,\{\alpha,\beta\})$-exit paths.  Their
outside ends $a_1,a_2,a_3$ and tree ends $b_1,b_2,b_3$ are pairwise
distinct, with $b_1\prec_T b_2\prec_T b_3$.  Retain the exact
$\beta$-colored boundary-edge identity $f_j$ of each path incident with
$b_j$.  Suppose $b_3\in V(A)$.

Assume that a selected color $\theta\ne\alpha,\beta$ is missing and
singleton-closed on $A$ and singleton-closed on $T$.  Assume also the
explicit no-smaller-presentation consequence of the global minimum: for
every proper, exactly one-uncolored coloring $\sigma$ stable from
$\varphi$ on the final ladder tree, arbitrary replacement input,
corresponding same-rung maximum ETT $T'$, good hierarchy with fewer than
$d$ dividers, its own large estimate, and its separate terminal
minus-closed certificate, there do not exist two support-disjoint exact
exit legs with distinct outside ends for a missing singleton-closed first
color.  Assume additionally that $d=0$ and that the terminal extension
kind is not PE.  Then the displayed Subcase 2.2 data are contradictory.
\end{lemma}
\begin{proof}
Write $n$ for the rung number, $T_i$ for the exact tree at series slot
$i$ (so $T_n$ is the last ladder tree), $D_i$ for the exact cumulative
stored-color set at slot $i$, and $\varphi_i$ for the exact coloring at
that slot.  Thus $\varphi=\varphi_n$ by the
generating-coloring field of \noderef{ExtendedTashkinovTree}.  We first
make the three exit legs used in every branch completely explicit.

The starred tree $A$ is an ordered initial segment of $T$.  Unfold
\noderef{TashkinovHierarchy}.  If $d>0$, then
\noderef{HierarchyStarredStage} identifies $A$ with the $d$-th hierarchy
stage, and repeated use of the nestedness field of the hierarchy gives
the required initial-segment map to its terminal stage $T$.  If $d=0$,
the same conclusion follows from the core-inside-first field and then the
nestedness and terminal-stage fields of the hierarchy.  All these are
entrywise statements by \noderef{TreeSequenceInitialSegment}.  Let
$\iota$ be the order-preserving inclusion of the vertex positions of $A$
into those of $T$ supplied by this initial-segment relation.  Unfolding
\noderef{TreeVertexOrder} in $b_1\prec_T b_2$ gives positions $i_1<i_2$
with $T(i_1)=b_1$ and $T(i_2)=b_2$; unfolding it in
$b_2\prec_T b_3$ gives positions $j_2<j_3$ with $T(j_2)=b_2$ and
$T(j_3)=b_3$.  The vertices in a \noderef{TreeSequence} are pairwise
distinct: applying its vertex-injectivity field to the two occurrences of
$b_2$ gives $i_2=j_2$.

Unfolding \noderef{TreeSequenceVertexSet} in the hypothesis
$b_3\in V(A)$ supplies a position $a_3$ of $A$ with $A(a_3)=b_3$.
Entrywise agreement gives $T(\iota(a_3))=A(a_3)=b_3$.  Apply the same
vertex-injectivity field, now to this occurrence and $T(j_3)=b_3$, to get
$\iota(a_3)=j_3$.  Hence $j_3$ lies in the initial block: numerically,
$\operatorname{val}(j_3)=\operatorname{val}(a_3)
\leq\operatorname{length}(A)$.  Since
$i_1<i_2=j_2<j_3$, the numerical positions $i_1$ and $j_2$ also define
positions $a_1,a_2$ of $A$.  Their images under $\iota$ are respectively
$i_1,j_2$ (the inclusions preserve numerical values), so entrywise
agreement gives $A(a_1)=b_1$ and $A(a_2)=b_2$.  Thus
$b_1,b_2,b_3\in V(A)$ by \noderef{TreeSequenceVertexSet}.

Unfold each given \noderef{TreeExitPath}.  The edge-node immediately
preceding the terminal vertex-node $b_j$ represents an exact boundary
edge of $T$: its other adjacent vertex-node on the path is outside $T$,
because $b_j$ is the only vertex-node of the path in $V(T)$.  By
\noderef{TwoColorIncidenceGraph} that edge has color $\alpha$ or
$\beta$.  It cannot have color $\alpha$, since
\noderef{TreeColorSetClosed} says that the exact $\alpha$-colored
boundary set is empty.  Hence this last boundary edge is
$\beta$-colored.  The separately retained identity $f_j$ has the same
two properties by hypothesis; unfolding \noderef{BoundaryColorEdges}
and \noderef{BoundaryEdges} shows in particular that its end other than
$b_j$ lies outside $T$.  Since $b_j\in V(A)\subseteq V(T)$, the same
identity $f_j$ is also a $\beta$-colored boundary edge of $A$.

We next justify which $\alpha,\theta$ components may be exchanged.  If an
$\alpha$-edge node occurs on $P_j$, it is not incident with $T$:
otherwise one of its incident vertex-nodes on this alternating path
would be in $T$, hence would be $b_j$, and the edge would be an
$\alpha$-colored boundary edge, contrary to the preceding closedness.
More generally, its full \noderef{TwoColorComponent} for
$\alpha,\theta$ contains no vertex of $T$.  For if it did, take a
shortest incidence-graph walk from the external edge to a vertex-node in
$T$ and its first vertex-node in $T$.  The preceding exact edge identity
crosses the boundary of $T$ and, by
\noderef{TwoColorIncidenceGraph}, has color $\alpha$ or $\theta$.
This contradicts respectively the assumed singleton closedness of
$\alpha$ or of $\theta$ on $T$.  The identical first-crossing argument
shows that every $\alpha,\theta$ component containing an outside vertex
is wholly outside $T$.

Exchange $\alpha$ and $\theta$ on every full $\alpha,\theta$ component
outside $T$, and call the result $\nu$.  Finiteness makes this a finite
choice of roots, and the preceding paragraph verifies the outside-support
clause of \noderef{OutsideKempeRecoloring}.  By
\noderef{OutsideKempeRecoloringStable}, $\nu$ is proper, leaves exactly
$e_0$ uncolored, fixes every exact edge incident with $T$, preserves the
missing-color set at each vertex of $T$, and is
$(T,C,\varphi_n)$-stable for every chosen control set $C$.  We verify
rather than assume that this stability restricts to $T_n$.  By
\noderef{ExtendedTashkinovTree}, $T_n$ is an ordered initial segment of
$T$.  Unfold \noderef{TreeSequenceInitialSegment}.  If
$v\in V(T_n)$, then \noderef{TreeSequenceVertexSet} supplies a position
of $v$ in $T_n$; mapping that position to the equal entry of $T$ proves
$v\in V(T)$.  Thus $V(T_n)\subseteq V(T)$.

Now fix an edge identity $f$ incident with $T_n$.  Unfolding
\noderef{TreeIncidentEdge} supplies a vertex $v\in V(T_n)$ incident with
$f$.  The preceding inclusion puts the same $v$ in $V(T)$, so the same
witness proves that $f$ is incident with $T$.  Likewise, if a color $c$
is missing on $T_n$ under $\varphi_n$, then unfolding
\noderef{TreeMissingColors} supplies $v\in V(T_n)$ at which $c$ is
missing; the same vertex lies in $V(T)$, and the same missingness witness
proves $c\in\overline{\varphi_n}(T)$.  Consequently
\[
 \overline{\varphi_n}(T_n)\subseteq
 \overline{\varphi_n}(T).
\]
For the controlled-edge clause of \noderef{StableColoring}, an edge
incident with $T_n$ is therefore incident with $T$, and a controlling
color in $\overline{\varphi_n}(T_n)\cup C$ belongs to
$\overline{\varphi_n}(T)\cup C$; the stable-edge clause on $T$ fixes that
edge.  For its vertex clause, every vertex of $T_n$ is a vertex of $T$,
so the missing-set equality on $T$ applies unchanged.  These are exactly
the two clauses saying that $\nu$ is
$(T_n,C,\varphi_n)$-stable, for every $C$.

For $j=1,2,3$, form a walk $Q_j$ with the same ordered vertex- and
edge-nodes as $P_j$.  Each old $\beta$-edge remains $\beta$ because
$\beta\ne\alpha,\theta$.  Every old $\alpha$-edge lies in one of the
selected outside components just proved and becomes $\theta$.  Thus
every successive pair of nodes of $P_j$ is adjacent in
$I_G(\nu;\theta,\beta)$, so $Q_j$ is a walk there; it is the same simple
node list and has the same positive length and support as $P_j$.
Its outside end still misses one of its two colors.  Explicitly, if
$\beta$ was missing at $a_j$, then the definition
\noderef{MissingColors}, together with \noderef{PresentColors}, says
that no incident edge had color $\beta$; none acquires that color in an
$\alpha,\theta$ exchange, so $\beta$ remains missing.  If instead
$\alpha$ was missing, then no incident edge had color $\alpha$.  The
whole $\alpha,\theta$ component through the outside vertex $a_j$ was
selected, so the $\theta$-edges incident with $a_j$ become $\alpha$ and
the new $\theta$-edges are exactly its old $\alpha$-edges; hence
$\theta$ is missing there under $\nu$.  We already proved that all three
tree ends lie in $A$, while $a_j\notin V(T)$ implies
$a_j\notin V(A)$.  Finally, the only vertex-node of $Q_j$ in $A$ is
$b_j$, because the corresponding claim for $P_j$ excludes every other
vertex even from the larger tree $T$.  This checks every field of
\[
       \operatorname{TreeExitPath}(A,\nu,\theta,\beta,a_j,b_j,Q_j).
       \tag{1}
\]
The supports remain pairwise disjoint, the three outside ends remain
pairwise distinct, and no claim about separation of the three
transformed full components has been made.

Apply \noderef{HierarchyLEMStable} to the just constructed outside
recoloring.  It supplies a corresponding ETT $\mathcal T^\nu$ displayed
on $T$, a hierarchy base $B^\nu$, a hierarchy $H^\nu$, and a good
hierarchy on it.  The corresponding certificate retains the rung, all
ordered ladder trees, exact connecting-edge sets, stored sets, and kinds;
the same-hierarchy conclusion retains every ordered stage, divider,
stored-color set, and Gamma pair.  The result also supplies maximum
property and transports the given terminal minus-closed certificate.
The stability established above identifies the missing-color sets on
$T_n$ under $\nu$ and $\varphi_n$, so the large bound is unchanged.

By the zero-divider hypothesis, $d=0$.  By \noderef{HierarchyStarredStage}, $A$ is the
actual hierarchy core.  We treat the two possible terminal constructions
separately.

The terminal extension is not PE by hypothesis.  The nonparallel field
of \noderef{HierarchyBase} identifies $A=T_n$.  Positivity permits us to
inspect the actual last certificate
$\mathcal T.\mathit{series.certified}(n)$ in
\noderef{TashkinovSeries}.  It cannot be the initial constructor.  The
PE constructor is excluded, so \noderef{TashkinovSeriesStep} leaves
exactly two cases.  In the RE case it gives the exact equation
$\varphi_n=\varphi_{n-1}$ and says $T_n$ is closed under
$\varphi_{n-1}$.  In the SE case it gives a
\noderef{MaximumDefectiveChoice} whose selected coloring is exactly
$\varphi_n$, whose selected-stability field says that $\varphi_n$ is
$(T_n,D_{n-1},\varphi_{n-1})$-stable, and again says that $T_n$ is closed
under $\varphi_{n-1}$.

We claim that every color missing somewhere on $T_n$ under $\nu$ is
singleton-closed there under $\nu$.  Fix such a color $c'$ and a witness
$y\in V(T_n)$.  Suppose an exact identity $g$ belonged to its
$\nu$-colored boundary.  By \noderef{BoundaryColorEdges}, $g$ is a
boundary edge and $\nu(g)=c'$.  The first outside recoloring is
$(T_n,D_{n-1},\varphi_n)$-stable.  Reverse it using \noderef{sc2}; its
two clauses are those of \noderef{StableColoring}: the vertex clause
transfers the missingness of $c'$ at $y$ to $\varphi_n$, and the edge
clause, now applied to the incident boundary identity $g$, gives
$\varphi_n(g)=\nu(g)=c'$.  In the RE case substitute the exact equality
$\varphi_n=\varphi_{n-1}$.  In the SE case reverse the selected stability
using \noderef{sc2}; the common missingness at $y$ and its edge clause
give $\varphi_{n-1}(g)=c'$.  In either case
\noderef{ColoredVertexSetClosed}, applied to the same exact boundary
identity $g$, the inside vertex $y$, and the color $c'$ missing at $y$,
is contradicted.  Thus the $c'$-colored boundary is empty, proving the
claim from the definition \noderef{TreeColorSetClosed}.

Choose any $x\in V(T_n)$.  The initial constructor of the certified
series supplies \noderef{KTriple}, hence $k\geq\Delta(G)+1$.
Properness means that the map from colors present at $x$ to incident
edges is injective: two distinct incident identities cannot both have
the same color.  Thus \noderef{PresentColors}, \noderef{Degree}, and
\noderef{MaximumDegree} give at most
$d_G(x)\leq\Delta(G)$ present colors.  At least one color $c$ is
therefore in \noderef{MissingColors}$(G,\nu,x)$.  It belongs to
$\overline\nu(T_n)$ by \noderef{TreeMissingColors}, and the preceding
claim makes $T_n$ singleton-$c$-closed.

Before making another recoloring, note that $c\ne\beta$.  The retained
identity $f_3$ is a $\beta$-colored boundary edge of $A=T_n$ under
$\varphi_n$, and the first outside recoloring fixes this exact incident
edge, so it is such a boundary edge under $\nu$ as well.  Equality
$c=\beta$ would contradict the singleton-$c$-closedness just proved.
The same boundary identity also proves that $\beta$ is present at $x$.
Indeed, if it were missing at $x$, the preceding claim, applied to
$c'=\beta$, would make the $\beta$-colored boundary of $T_n$ empty,
contrary to $f_3$.  By the complementary definitions
\noderef{MissingColors} and \noderef{PresentColors}, there is therefore
an exact $\beta$-colored edge incident with $x$.

If $\theta=c$, put $\rho=\nu$ and retain the walks $Q_j$.  Otherwise,
exchange $\theta$ and $c$ on every full component outside $T_n$ and call
the result $\rho$.  The first recoloring fixed the boundary of $T_n$, so
the assumed singleton-$\theta$-closedness persists under $\nu$;
together with singleton-$c$-closedness, the same first-crossing argument
used above proves that these are genuine outside components.  Hence
\noderef{OutsideKempeRecoloringStable} gives properness, exact
one-uncoloredness, unchanged missing sets on $T_n$, and
$(T_n,C,\nu)$-stability for every control set $C$.  The identity case has
the same conclusions by reflexivity in \noderef{sc2}.

Exactly as in the construction of (1), replace each $\theta$-edge of
$Q_j$ by its new label $c$ and retain each exact $\beta$-edge, for
$j=1,2,3$.  The inequalities $\theta\ne\beta$ and $c\ne\beta$ ensure
that $\beta$ is never exchanged.  If the old outside end missed
$\beta$, it still does; if it missed $\theta$, all incident
$\theta,c$ edges lie in its selected outside component and its new
$c$-edges are precisely its old $\theta$-edges, so it misses $c$.
The tree ends and the unique-tree-node condition are unchanged.  We thus
obtain three exact $(T_n,\rho,\{c,\beta\})$-exit paths
$L_1,L_2,L_3$ with the same respective supports as $P_1,P_2,P_3$.
They remain pairwise disjoint and their outside ends remain pairwise
distinct.  The fixed boundary identity $f_3$ also shows directly that
$\beta$ remains present at $x$ under $\rho$, while stability preserves
the missingness of $c$ there.

We next certify the full path components that will contradict
interchangeability.  In \noderef{TwoColorIncidenceGraph}, the
vertex-node $x$ has no $c$-colored neighbor because $c$ is missing, and
it has at least one $\beta$-colored neighbor because $\beta$ is present.
Unfolding \noderef{IsProperEdgeColoring}, properness allows at most one
edge of either color at a vertex, so $x$ has degree exactly one.  Each
outside end $a_j$ also has degree exactly one: the nontrivial exit path
$L_j$ supplies an incident edge of one of the two colors, its
missing-color clause excludes the other color, and properness excludes a
second edge of the present color.  In particular,
$x,a_1,a_2,a_3$ are four distinct degree-one vertex-nodes, since
$x\in V(T_n)$ whereas the $a_j$ lie outside it and the $a_j$ are
pairwise distinct.

Let $X$ be the full \noderef{TwoColorComponent} through $x$, and let
$Y_j$ be the full component through $a_j$.  The degree-one witnesses
just proved verify \noderef{TwoColorPathComponent} for each of
$X,Y_1,Y_2,Y_3$ (and $c\ne\beta$ supplies its distinct-color field).
The component $X$ meets $T_n$ at $x$; connectivity along $L_j$ puts
$b_j\in V(T_n)$ in $Y_j$.  At least one of $Y_1,Y_2,Y_3$ is different
from $X$.  Otherwise one component would contain the four distinct
degree-one nodes $x,a_1,a_2,a_3$.  To see the contradiction without
assuming separation of the transformed legs, use
\noderef{TwoColorComponentSpanningPath}: a simple path spans that full
component.  Any ambient degree-one node in its support must be an
endpoint of the spanning path, since an internal node has its two
distinct predecessor and successor as neighbors.  A path has only two
endpoints.  This also accords with the global degree bound
\noderef{TwoColorIncidenceDegreeAtMostTwo}.  Choose $j$ with $X\ne Y_j$.

We verify the exact stable-coloring premise needed to make these two path
components equal.  The second outside recoloring gives $\rho$
$(T_n,D_{n-1},\nu)$-stable, and the first gives $\nu$
$(T_n,D_{n-1},\varphi_n)$-stable.  In the RE case substitute
$\varphi_n=\varphi_{n-1}$; in the SE case append the selected-stability
field, which is exactly
$(T_n,D_{n-1},\varphi_{n-1})$-stability of $\varphi_n$.
Transitivity in \noderef{sc2} therefore gives
\[
        \rho\text{ is }(T_n,D_{n-1},\varphi_{n-1})\text{-stable}.
        \tag{2}
\]
All intermediate colorings are proper and leave precisely $e_0$
uncolored, by the two outside-recoloring results and the corresponding
certified constructors.

If $n>1$, put $m=n-1$.  Applying
\noderef{CertifiedETTPredecessor} to the original ETT gives an actual
positive-rung predecessor of rung $m$, with maximum property, displayed
terminal tree exactly $T_n$, and generating coloring exactly
$\varphi_{n-1}$.  Its statement identifies every retained tree slot and,
at every retained data index, the coloring and raw stored-color slot with
the corresponding original entries.  We now derive the equality of the
computed control sets, which is not merely a raw-slot equality.  Write
$S^-_h,\varphi^-_h,T^-_h$ for the predecessor's raw stored, coloring,
and tree entries at a retained numerical index $h$.

At the predecessor's last data index $m>0$, unfold
\noderef{TashkinovStoredColors}.  Its computed set is
\[
 \left(\bigcup_{0\leq h\leq m}S^-_h\right)
 \setminus
 \overline{\varphi^-_m}(T^-_m).                         \tag{3}
\]
The canonical retained-index inclusion preserves numerical values, and
$m=n-1$; hence its indices $h\leq m$ are in bijection with exactly the
original indices $h\leq n-1$.  The raw stored-slot equalities from
\noderef{CertifiedETTPredecessor} identify every $S^-_h$ in (3) with
$S_h$, so they identify the two finite unions.  Its coloring equality at
$m$ gives $\varphi^-_m=\varphi_{n-1}$, and its tree-slot equality at the
same numerical index gives $T^-_m=T_{n-1}$.  Unfolding
\noderef{TreeMissingColors}, these two equalities identify the set
deleted in (3) with
$\overline{\varphi_{n-1}}(T_{n-1})$.  A second use of
\noderef{TashkinovStoredColors} therefore turns (3) into the original
computed set $D_{n-1}$.  Notice separately that the predecessor tree
slot after this control index is slot $m+1=n$; the retained tree-slot
equality identifies it with the original $T_n$, exactly the tree used in
lower assertion (ii).  Thus (2), properness, exact one-uncoloredness, and
the color $c$ missing at $x$ satisfy that assertion for this predecessor.

If $n=1$, the relevant original data index is $n-1=0$.  The zero branch
in the definition \noderef{TashkinovStoredColors} is taken before any
finite union or missing-color deletion and gives
$D_{n-1}=D_0=\varnothing$ explicitly.  The initial certificate in
\noderef{TashkinovSeries}, through the initial constructor of
\noderef{TashkinovSeriesStep}, has next tree $T_1=T_n$ and coloring
$\varphi_0=\varphi_{n-1}$.  Consequently (2) says precisely that $\rho$
is $(T_1,D_0,\varphi_0)$-stable, and its properness and exact
one-uncoloredness are the remaining hypotheses of
\noderef{TechnicalTheoremBase}.  In both cases the result says that
$c,\beta$ are $T_n$-interchangeable.  Unfold
\noderef{TreeInterchangeable} and apply it to the two certified path
components $X$ and $Y_j$, each of which meets $T_n$.  It gives
$X=Y_j$, contrary to their choice.  No zero-rung ETT has been used.

\end{proof}
